% TeXFrog tutorial source: IND-CPA proof for PRF-based symmetric encryption.
%
% Scheme: Enc(k, m) = (r, PRF(k, r) xor m),  where r is sampled fresh each call.
%
% Tag syntax: %:tags: label1,label2-label3
% - Lines with no %:tags: comment appear in every game/reduction.
% - Lines tagged with one or more labels appear only in those games/reductions.
% - Ranges like G0-G2 are resolved by POSITION in the 'games:' list, not alphabetically.
% - Variant lines for the same "slot" must be consecutive in this file.

\begin{pcvstack}[boxed]

%%% -- Challenge experiment procedure header (one per game/reduction) -------------

    \procedure[linenumbering]{Game $\tfgamename{G0} = \INDCPA_\Enc^\Adversary.\mathsf{Real}()$}{ %:tags: G0
    \procedure[linenumbering]{Game $\tfgamename{G1}$}{ %:tags: G1
    \procedure[linenumbering]{Game $\tfgamename{G2} = \INDCPA_\Enc^\Adversary.\mathsf{Ideal}()$}{ %:tags: G2
    \procedure[linenumbering]{Reduction $\Bdversary_1^{\OPRF}$}{ %:tags: Red1

%%% -- Challenge experiment body -------------------------------------------------

        % PRF key: G0, G1, G2 only (Red1 uses an external PRF oracle, not a local key)
        k \getsr \{0,1\}^\lambda \\ %:tags: G0-G2
        % Challenge bit (all games)
        b \getsr \{0,1\} \\
        % Run adversary with access to the LR encryption oracle (all games)
        b' \gets \Adversary^{\mathsf{LR}}() \\
        \pcreturn (b' = b)
    }

%%% -- Left-or-Right encryption oracle -------------------------------------------

    \procedure[linenumbering]{$\mathsf{LR}(m_0, m_1)$}{

        % Fresh nonce for every query (all games)
        r \getsr \{0,1\}^\lambda \\
        % G0: compute the PRF output using the secret key
        y \gets \mathrm{PRF}(k, r) \\ %:tags: G0
        % G1: use a truly random value instead of the PRF output
        y \getsr \{0,1\}^\lambda \\ %:tags: G1
        % Red1: query the external PRF oracle (no local key)
        y \gets \OPRF(r) \\ %:tags: Red1
        % G0, G1, Red1: XOR with the chosen message
        c \gets y \oplus m_b \\ %:tags: G0,G1,Red1
        % G2: ciphertext is uniformly random, independent of m_b
        c \getsr \{0,1\}^\lambda \\ %:tags: G2
        \pcreturn (r, c)
    }

\end{pcvstack}
